% 参数	默认值	控制内容
% body-width	0.86\textwidth	整个红框内的信息区域宽度
% line-height	24bp	正文信息的基线间距
% row-sep	2mm	委员、秘书、时间、地点之间的附加间距
% role-width	4\ccwd	右侧主席/委员的预留宽度
% role-gap	0.5\ccwd	下划线终点与角色之间的距离
% name-align	column	姓名采用统一中心线还是逐条居中, line
% name-position	0.63	统一姓名中心线距信息区域左端的比例，仅 column 模式使用



% 模式	行为
% name-align=column（默认）	所有姓名的中心在同一竖直线上，适合你提供的 Word 样式
% name-align=line	每个姓名严格位于自己所在横线的正中，但不同姓名可能不在同一竖直线上

%
% name-align = column,
% name-position = 0.65,
% 表示所有姓名的中心位置设在信息区域宽度的 65% 处，姓名整体比默认位置略靠右。所有名字仍然处在同一竖直线上。




\begin{committee}[
  name = {硕士学位论文答辩委员会名单},
  body-width = 0.8\textwidth,
  line-height = 34bp,
  row-sep = 0mm,
  name-align = column,
  name-position = 0.63,
  role-gap = 0.5\ccwd
]

\committeemember{XXXX大学XXX教授}{xxx}{主席}
\committeemember{XXXX大学XXX教授}{xxx}{委员}
\committeemember{贵州大学教授/博导}{xxx}{委员}
\committeemember{贵州大学教授/博导}{xxx}{委员}

\committeesecretary[
  before-name = 1\ccwd,
  before-title = 1\ccwd
]{XXX}{副教授/博士研究生}

\committeedate{XXXX年XX月XX日}
\committeeplace{贵州大学西校区崇厚楼XXXX}

\end{committee}



% 也可以导入 Word 版转的 PDF 文件
% \begin{committee}[file=figures/committee.pdf]
% \end{committee}